\section{All square-zero extensions of the localized integers}
\label{sec:arbitrary-idealizations}

Fix a prime $p$ and put $O=\mathbb Z_{(p)}$.  For an arbitrary
$O$-module $M$, without a finite-generation hypothesis, write
\[
 O\ltimes M=O\oplus M,\qquad
 (a,m)(b,n)=(ab,an+bm).
\]
This is a local characteristic-zero ring with square-zero ideal $M$
and residue field $\mathbb F_p$.

\begin{theorem}\label{thm:arbitrary-idealizations}
For every $O$-module $M$, the polynomial map generated by the trivial
symmetric-group classes is an isomorphism over $O\ltimes M$ if and
only if $M$ is a $\mathbb Q$-vector space.  In every negative case,
surjectivity already fails in degree $p$.

If $M$ has nonzero torsion, then already in symmetric-group degree
$p$ the kernel of augmentation
\[
 K_0\bigl(\operatorname{Lat}_{O\ltimes M}(S_p)\bigr)
 \longrightarrow K_0\bigl(\operatorname{Lat}_{O}(S_p)\bigr)
\]
contains a direct summand isomorphic to a countably infinite direct
sum of copies of $\mathbb Z$.
In particular this assertion includes the divisible torsion module
$M=\mathbb Q/O$.
\end{theorem}

\subsection{An integer detector for the divisible torsion extension}

Put $D=\mathbb Q/O$ and $B=O\ltimes D$.  Fix
$H=C_p=\langle s\rangle$.  If $E$ is a finite $B$-free $B[H]$-module,
put $E_0=E/DE$.  Freeness gives a canonical identification
$DE=D\otimes_OE_0$ and canonical exact sequences
\begin{equation}\label{eq:divisible-two-extensions}
 0\longrightarrow D\otimes_OE_0\longrightarrow E\longrightarrow E_0
 \longrightarrow0,
 \qquad
 0\longrightarrow E_0\longrightarrow\mathbb Q\otimes_OE_0
 \longrightarrow D\otimes_OE_0\longrightarrow0.
\end{equation}
Use the two-periodic cochain complex with even differential $S-I$
and odd differential $I+S+\cdots+S^{p-1}$.  Write its cohomology as
$H^j$, with the indicated identification $H^{j+2}=H^j$.
The two groups $H^0(E_0),H^1(E_0)$ are finite-dimensional
$\mathbb F_p$-spaces: they are finitely generated over $O$ and are
killed by $p$, using the norm on invariants and the identity
\[
 (S-I)\sum_{i=0}^{p-2}(p-1-i)S^i
 =I+S+\cdots+S^{p-1}-pI.
\]

Let $\delta_E^j:H^j(E_0)\to H^{j+1}(D\otimes E_0)$ and
$\beta_E^j:H^j(D\otimes E_0)\to H^{j+1}(E_0)$ be the boundaries
of~\eqref{eq:divisible-two-extensions}.  Define
\begin{equation}\label{eq:divisible-u}
 u_E^j=\beta_E^{j+1}\delta_E^j:
 H^j(E_0)\longrightarrow H^{j+2}(E_0)=H^j(E_0).
\end{equation}
Both boundary maps are canonical.  The second is an isomorphism
because the rational middle complex is acyclic, although only its
naturality is needed below.

We check compatibility with admissible exact sequences explicitly.
Choose a $B$-linear splitting of $0\to E'\to E\to E''\to0$.
The cochain differential has upper-triangular block form; denote its
constant diagonal blocks by $d_0',d_0''$, its constant off-diagonal
block by $U_0$, and its $D$-valued blocks by $d_1',d_1'',U_1$.
The square-zero part of $d^{j+1}d^j=0$ is
\[
 d_1'{}^{j+1}U_0^j+U_0^{j+1}d_1''{}^j
 =-d_0'{}^{j+1}U_1^j-U_1^{j+1}d_0''{}^j.
\]
On an integral cocycle the last term vanishes and the first term
on the right is a boundary in $D\otimes E_0'$.  Thus, writing
$\partial_0,\partial_D$ for the connecting maps of the reduced
and $D$-coefficient sequences, we obtain
\begin{equation}\label{eq:divisible-delta-compatibility}
 \delta_{E'}\partial_0=-\partial_D\delta_{E''}.
\end{equation}
For the second boundary, represent a $D$-coefficient cocycle $z$
by a rational lift $\widetilde z$.  Its differential is an integral
cocycle $w$, which represents $\beta_{E''}[z]$.  The constant block
identity gives
\[
 d_0'{}^{j+1}U_0^j\widetilde z
 =-U_0^{j+1}d_0''{}^j\widetilde z=-U_0^{j+1}w.
\]
Consequently
\begin{equation}\label{eq:divisible-beta-compatibility}
 \beta_{E'}\partial_D=-\partial_0\beta_{E''}.
\end{equation}
The two signs cancel in the composite, so
$u_{E'}\partial_0=\partial_0u_{E''}$.  All other maps commute with
$u$ by functoriality.

For a monic irreducible $f\in\mathbb F_p[t]$, $f\ne t$, define
\begin{equation}\label{eq:divisible-euler-detector}
 \chi_f(E)=
 \ell_{\mathbb F_p[t]}\bigl(H^0(E_0)_{f\text{-primary}}\bigr)
 -\ell_{\mathbb F_p[t]}\bigl(H^1(E_0)_{f\text{-primary}}\bigr),
 \qquad t=u_E.
\end{equation}
Taking primary parts preserves exactness: a common annihilator
$f^a q$, with $(f,q)=1$, gives a common primary projector by Bezout's
identity.  The six-term periodic exact cycle therefore proves that
$\chi_f$ is an integer-valued additive invariant.  It vanishes on
every lattice extended from $O$, because the first sequence
in~\eqref{eq:divisible-two-extensions} then splits equivariantly and
$\delta_E=0$.  Each class has only finitely many nonzero $\chi_f$.

Let $\eta=1/p+O\in D$.  On the rank-one $B[H]$-module with
$s=1+a\eta$, $a\in\mathbb F_p$, the integral reduction is trivial,
with $H^0=\mathbb F_p$ and $H^1=0$.  The first boundary sends $1$ to
$a\eta$, and the second lifts it to $a/p$ and applies the norm $p$.
Thus $u=a$ on the even group.  More generally, for a matrix $C$ over
$\mathbb F_p$, the action $I+\eta C$ has $u=C$ on its even Tate
group and has zero odd integral Tate group.

\subsection{Separating arbitrary nonzero torsion}

Suppose $M$ has nonzero torsion.  Since $O$ is a discrete valuation
ring, some nonzero $m\in M$ satisfies $pm=0$.  Define
$Om\to D$ by $m\mapsto\eta$.  This extends to an $O$-linear map
\begin{equation}\label{eq:torsion-separating-functional}
 \lambda:M\longrightarrow D.
\end{equation}
Here is an elementary proof of the required extension property.
Use Zorn's lemma on partial extensions.  If a partial map is
defined on $N\subset M$ and $x\notin N$, the ideal
$\{a\in O:ax\in N\}$ is principal, say $(d)$.  If $d=0$, choose
the image of $x$ arbitrarily.  If $d\ne0$, divisibility of $D$
provides $y\in D$ with $dy=\lambda(dx)$.  The formula
$\lambda(n+ax)=\lambda(n)+ay$ is then well-defined on $N+Ox$.
Thus a maximal partial extension has domain $M$.

Put $A=O\ltimes M$.  The map~\eqref{eq:torsion-separating-functional}
induces a coefficient homomorphism $A\to B$.  Let $L=A$ with
$s$ acting as $1+m$, and set
\[
 N=\operatorname{Ind}_{H}^{S_p}L,\qquad
 N_0=A\otimes_O\operatorname{Ind}_{H}^{S_p}O.
\]
This is a representation since $(1+m)^p=1$.
Apply base change to $B$, restriction to $H$, and then
$\chi_{t-1}$.  In the Mackey decomposition, a double coset has
intersection $H$ precisely when its representative normalizes
$H$; all other intersections are trivial.  The normalizer quotient
is $\mathbb F_p^\times$, with each automorphism represented once,
and the corresponding rank-one twists have $u=a$ for
$a\in\mathbb F_p^\times$.  The other summands are induced from
the trivial subgroup and have zero reduced Tate cohomology.
Consequently the composite integer invariant takes
\[
 [N]-[N_0]\quad\hbox{to}\quad1.
\]
It vanishes on all Young permutation classes.  This proves failure
of generation in degree $p$, and also proves that the displayed
class spans a split copy of $\mathbb Z$ in the augmentation kernel.

For odd $p$ this class has zero trace at every permutation.  At a
$p$-cycle, the two induced traces differ by
$\sum_{a=1}^{p-1}am=p(p-1)m/2=0$, since $pm=0$.
At the identity the ranks agree; every other nonidentity cycle type
is not conjugate into $H$, so both induced traces vanish.
For the smallest example $M=\mathbb F_p\epsilon$, the cyclic
line has action $1+\epsilon$.  Its induced trace is therefore the
same as the constant induced line when $p$ is odd, despite the
infinite order detected above.


For the countably infinite direct summand, choose strictly
increasing positive integers $r_i$ coprime to $p-1$, and monic
irreducibles $f_i\in\mathbb F_p[t]$ of degree $r_i$ with nonzero
constant term.  Such irreducibles exist in every positive degree.
Let $C_i$ be a companion matrix, and let $L_i=A^{r_i}$ with action
$I+mC_i$.  Put
\[
 x_i=[\operatorname{Ind}_{H}^{S_p}L_i]-r_i[N_0].
\]
Twisting by $s\mapsto s^a$ replaces $C_i$ by $aC_i$ and its
polynomial by $a^{r_i}f_i(t/a)$.  Equality with $f_i$ forces
$a^{r_i}=1$ by comparison of constant terms, hence $a=1$ because
$(r_i,p-1)=1$.  Different chosen degrees have disjoint scaling
orbits.  The induced invariant family therefore satisfies
\[
 \chi_{f_j}\bigl(\operatorname{Res}_{H}^{S_p}(x_i\otimes_A B)\bigr)
 =\delta_{ij}.
\]
It has finite support on each class and so maps to
$\bigoplus_i\mathbb Z$.  Sending its standard basis vectors to
$x_i$ is a section.  All $x_i$ reduce to zero over $O$, proving the
direct-summand assertion in the theorem.

\subsection{The torsion-free alternatives}

\begin{lemma}\label{lem:padic-separation}
Let $O=\Z_{(p)}$, let $M$ be a torsion-free $O$-module, and let
$m\in M\setminus pM$.  There is an $O$-linear map
\[
 \varphi:M\longrightarrow\Z_p,\qquad \varphi(m)=1.
\]
\end{lemma}

\begin{proof}
Choose an $\F_p$-basis $(\overline m_s)_{s\in S}$ of $M/pM$
containing $\overline m$, and choose representatives $m_s\in M$, with
$m_{s_0}=m$.  These representatives are $O$-linearly independent.
Indeed, divide the coefficients in a nonzero finite relation by their
smallest common power of $p$.  Torsion-freeness allows cancellation of
that power in $M$, and reduction modulo $p$ would then give a nonzero
relation among the chosen basis vectors.  Consequently
\[
 F=\sum_{s\in S}O m_s=\bigoplus_{s\in S}O m_s
\]
is a free submodule of $M$, and $F/pF\to M/pM$ is an isomorphism.

For every $n\geq1$ the inclusion induces an isomorphism
\begin{equation}\label{eq:idealization-completion}
 F/p^nF\xrightarrow{\ \sim\ }M/p^nM.
\end{equation}
For surjectivity, successive approximation modulo $p$ writes an arbitrary
$x\in M$ in the form
\[
 x=f_0+pf_1+\cdots+p^{n-1}f_{n-1}+p^n x_n,
 \qquad f_j\in F,\quad x_n\in M.
\]
For injectivity, if $f\in F\cap p^nM$, injectivity modulo $p$ first
gives $f=pf_1$ with $f_1\in F$.  Cancellation of $p$ in $M$ gives
$f_1\in p^{n-1}M$; iteration yields $f\in p^nF$.

The coordinate projection $F\to O$ at $s_0$ induces compatible maps
$F/p^nF\to O/p^nO$.  Transporting them across
\eqref{eq:idealization-completion} and taking their inverse limit gives
\[
 M\longrightarrow\varprojlim_n O/p^nO=\Z_p.
\]
It sends $m$ to $1$.  Only finite linear combinations occur at every
stage, so the proof imposes no restriction on the cardinality of $S$.
\end{proof}


We include the reduction to the flat dual-number obstruction for
completeness.  Suppose $M$ is torsion-free and $M\ne pM$.  Choose $m\notin pM$.  Lemma~\ref{lem:padic-separation} gives
an $O$-linear map $\ell:M\to\mathbb Z_p$ with $\ell(m)=1$.
Hence there are coefficient maps
\[
 O[\epsilon]/(\epsilon^2)\longrightarrow O\ltimes M
 \longrightarrow\mathbb Z_p[\epsilon]/(\epsilon^2),
 \qquad \epsilon\longmapsto m\longmapsto\epsilon.
\]
The explicit rank-$p$ cyclic lattice and its induced degree-$p$
counterexample from the flat dual-number theorem can be extended
along these maps.  Its parity detector is still nonzero over
$\mathbb Z_p$: the two-periodic construction uses finite integral
Tate groups over the residue field, and the displayed boundary
matrices are unchanged by completion.  Equivalently the proof of
Theorem~\ref{thm:dual-numbers-all-primes} applies verbatim to the
DVR $\mathbb Z_p$.  A Young expression over $O\ltimes M$ would
give one after the second base change, a contradiction.

It remains to prove the positive alternative.  A torsion-free
$O$-module with $M=pM$ is exactly a $\mathbb Q$-vector space,
because multiplication by $p$ is then bijective.  Suppose $M$ is
such a vector space and let $G$ be any finite group.  In an
$A$-basis, every representation can be written
\[
 \rho(g)=(I+c_g)\rho_0(g),\qquad c_g\in\operatorname{Mat}_r(M),
\]
where $\rho_0$ is its reduction over $O$.  Conjugation by $\rho_0$
gives the action on the matrix space, and the representation
identity is the cocycle equation $c_{gh}=c_g+g\cdot c_h$.
Define
\[
 X=\frac1{|G|}\sum_{g\in G}c_g.
\]
Then $g\cdot X=X-c_g$.  Since the ideal is square-zero,
conjugation by $I-X$ gives
\[
 (I-X)\rho(g)(I+X)=\rho_0(g).
\]
Thus every lattice is isomorphic to a scalar extension from $O$.
The polynomial theorem over $O$, together with the augmentation
retraction, proves the polynomial isomorphism over $A$.
The torsion and torsion-free cases exhaust all modules, completing
the proof of Theorem~\ref{thm:arbitrary-idealizations}.

\begin{corollary}\label{cor:all-square-zero-extensions}
Let $I^2=0$ in a commutative ring $A$ and suppose $A/I\simeq
\mathbb Z_{(p)}$.  Then the polynomial map over $A$ is an
isomorphism if and only if $I$ is a $\mathbb Q$-vector space under
its natural $\mathbb Z_{(p)}$-module structure.
If $I$ is finitely generated, this holds if and only if $I=0$.
In particular, for every $n\ge1$,
\[
 \mathbb Z_{(p)}[\epsilon]/(p^n\epsilon,\epsilon^2)
\]
has a split countably generated free abelian summand in the
degree-$p$ augmentation kernel.
\end{corollary}

\begin{proof}
Every integer prime to $p$ is a unit modulo $I$ and therefore is
a unit in $A$.  The localization property supplies a unique
section $\mathbb Z_{(p)}\to A$ of the quotient map.  Thus
$A\simeq O\ltimes I$, and the theorem applies.  If $I$ is finitely
generated and $I=pI$, Nakayama's lemma gives $I=0$.
For the displayed cyclic torsion ideal, one may take
$m=p^{n-1}\epsilon$ in the construction; it has order $p$.
\end{proof}

No classification of nonsplit higher nilpotent extensions or of
arbitrary singular local rings is asserted here.
